\section{Energy compression and polynomial IID learning}

The IID argument uses only boundedness and monotonicity of the activation.
It does not require that all low-energy output functions be affine
thresholds. A bounded monotone neuron is an average of threshold neurons;
sampling from that average gives a finite class containing every
low-energy full-margin sign function.

\begin{lemma}[Compression into a majority of thresholds]\label{lem:compression}
Let $\phi:\R\to[-M,M]$ be nondecreasing, where $M>0$, and let $E\ge0$.
Consider all shallow networks on $\{-1,1\}^d$, over all finite widths,
whose energy~\eqref{eq:energy} is at most $E$ and whose outputs have
absolute margin at least $1/2$ everywhere. Every resulting Boolean
function is an unweighted majority of
\[
 m_* = \left\lceil256(E+1)M^2(d+1)\right\rceil
\]
affine threshold functions. Consequently the class $\mathcal G_E$ of
these Boolean functions satisfies
\begin{equation}\label{eq:compression-count}
 \log|\mathcal G_E|\le m_*(d+1)^2\log2.
\end{equation}
Ties in the majority may be resolved by any fixed convention.
\end{lemma}
\begin{proof}
Write $F(x)=\sum_j v_j\phi(u_j\cdot\bar x)$ and
$B=\sum_j|v_j|$. Cauchy--Schwarz gives
\[
 B^2\le N\sum_jv_j^2\le E.
\]
The full margin implies $B>0$.
For a fixed neuron, choose $U$ uniformly from $[0,1]$ and define
\[
 h_U(x)=
 \begin{cases}
 +1,&U\le[\phi(u_j\cdot\bar x)+M]/(2M),\\
 -1,&U>[\phi(u_j\cdot\bar x)+M]/(2M).
 \end{cases}
\]
Then $\E_U h_U(x)=\phi(u_j\cdot\bar x)/M$.
Monotonicity means that every such upper level set is a threshold in
$u_j\cdot\bar x$. On the finite cube, almost every $U$ avoids all the
finitely many activation values used here. Its truth table is therefore
a strict affine threshold function, including a possible constant table.
Indeed one can place an affine threshold between the largest negative
score and the smallest positive score. This argument allows jumps in
$\phi$ and requires no continuity or parameter rounding.

Now select neuron $j$ with probability $|v_j|/B$, sample its level $U$,
and multiply its threshold table by $\sgn(v_j)$. Negating a strict
threshold table gives another strict threshold table. Thus there is a
probability distribution over affine threshold functions $h$ such that
\[
 F(x)=MB\,\E h(x).
\]
Take $m_*$ independent draws $h_1,\ldots,h_{m_*}$. Hoeffding's inequality
and a union bound over the cube give
\[
 \Pp\left(\max_x\left|
    \frac{MB}{m_*}\sum_{j=1}^{m_*}h_j(x)-F(x)
       \right|\ge\frac18\right)
 \le 2^{d+1}\exp\left(-\frac{m_*}{128M^2B^2}\right)<1.
\]
For the last inequality, use $B^2\le E$ and the definition of $m_*$;
when $E=0$ the full-margin class is empty.
There is therefore a realization with uniform error less than $1/8$.
Its sum has the same sign as $F$ at every input and has no zero value.
Multiplication by the positive scalar $MB/m_*$ does not affect that
sign, proving the majority representation.

Let $\mathcal H_d$ denote the affine threshold truth tables on the cube.
The $2^d$ input points determine hyperplanes in a $(d+1)$-dimensional
coefficient space. Adding one hyperplane at a time gives the usual
region bound
\[
 |\mathcal H_d|\le\sum_{j=0}^{d+1}\binom{2^d}{j}
 \le(2^d+1)^{d+1}\le2^{(d+1)^2}.
\]
There are at most $|\mathcal H_d|^{m_*}$ ordered threshold lists of
length $m_*$. This proves~\eqref{eq:compression-count}.
\end{proof}

\begin{theorem}[IID guarantee]\label{thm:iid}
Fix $\epsilon>0$, an affine threshold target $f$ on the cube, an input
distribution $\mathcal D$, and $\eta,\beta\in(0,1)$. Put
\[
 E_\epsilon=\epsilon^{-2},\qquad M_\epsilon=1+\epsilon,
 \qquad
 m_\epsilon=
 \left\lceil256(E_\epsilon+1)M_\epsilon^2(d+1)\right\rceil,
 \qquad \Lambda_\epsilon=m_\epsilon(d+1)^2\log2.
\]
If the number of IID training examples satisfies
\begin{equation}\label{eq:iid-sample}
 n\ge\frac4\eta
       \left[\Lambda_\epsilon+\log(1/\beta)\right],
\end{equation}
then, with sample probability at least $1-\beta$, the posterior
correct-label probability is greater than $2/3$ outside a set of
$\mathcal D$-mass at most $\eta$ in each of these regimes:
\begin{enumerate}
\item at every fixed width, for all sufficiently small $\sigma>0$;
\item under~\eqref{eq:scaling}, at every sufficiently large width for
any fixed $\gamma>1$ and $\tau>0$;
\item at every width and every $\gamma\ge1$ under~\eqref{eq:scaling},
for every sufficiently small fixed $\tau>0$.
\end{enumerate}
The thresholds can be chosen independently of the realized sample and
$\mathcal D$. At fixed $\epsilon$, the sample bound is
$O_\epsilon((d^3+\log(1/\beta))/\eta)$.
\end{theorem}
\begin{proof}
Choose any affine separator $a$ with $f(x)a\cdot\bar x\ge1$ on the
whole cube, and let $Q=\|a\|^2$. Cloning $u_j=ra$ and
$b_j=1/[2(\epsilon+\tanh r)]$ gives a full-target weak-margin witness
of energy
\[
 dQr^2+\frac1{4(\epsilon+\tanh r)^2}.
\]
Its derivative at $r=0$ is negative, so a sufficiently small positive
$r$ gives energy $e_f<1/(4\epsilon^2)<E_\epsilon$.
Every preactivation is nonzero. This upper bound needs no covariance
certificate for the target.
Apply Lemma~\ref{lem:box} to this witness and the event
$E_N\ge E_\epsilon$. The witness fits every target-consistent history,
so the resulting posterior bound is uniform over all such histories.
It tends to zero in the first two regimes and can be made at most
$\eta/12$ at every width in the third.

Lemma~\ref{lem:compression} gives a single class $\mathcal G_{E_\epsilon}$,
independent of learner width, containing every full-margin function
represented below that energy threshold. Its logarithmic size is at
most $\Lambda_\epsilon$. A consistency union bound and
\eqref{eq:iid-sample} imply, with probability at least $1-\beta$, that
every sample-consistent member of this class has population error at
most $\eta/4$. On this event, the posterior mean population error is
at most $\eta/4+\eta/12=\eta/3$.
An input with correct-label posterior probability at most $2/3$ has
wrong-label posterior probability at least $1/3$. Markov's inequality
therefore bounds the mass of these inputs by $\eta$, proving the claim.
\end{proof}
